% ID: SEL-D0020-OA-E03
% ESTADO: validada. S021, revisada con los criterios oficiales.
Los tres grupos tienen 220, 208 y $520-220-208=92$ usuarios.
Usan las duchas $220/5=44$ niños, $0.75\cdot208=156$ adultos menores
de 60 años y 23 del grupo restante. Son $44+156+23=223$ en total.

\begin{enumerate}[label=\alph*)]
\item Entre los 520 usuarios equiprobables, 223 se duchan.
La probabilidad es $\boxed{223/520}$.
\item La intersección pedida contiene a los 156 adultos menores
de 60 años que usan las duchas. Resulta $\boxed{156/520=0.30}$.
\item La condición reduce el universo a los 223 que se duchan.
De ellos 44 son niños, así que la probabilidad es $\boxed{44/223}$.
\end{enumerate}
